
\documentclass[a4paper,11pt]{article}
%%%%%%%%%%%%%%%%%%%%%%%%%%%%%%%%%%%%%%%%%%%%%%%%%%%%%%%%%%%%%%%%%%%%%%%%%%%%%%%%%%%%%%%%%%%%%%%%%%%%%%%%%%%%%%%%%%%%%%%%%%%%
\usepackage[brazilian]{babel}
\usepackage[T1]{fontenc}
\usepackage[dvips,dvipdf]{graphicx,epsfig}
\usepackage{subfigure}
\usepackage{url}
\usepackage{makeidx}
\usepackage[dvips,dvipdf]{colortbl}

\parindent15pt  \parskip0pt              
\setlength\voffset{-2.0cm}
\setlength\hoffset{-1.5cm}
\setlength\textwidth{16.0cm}
\setlength\textheight{24.5cm}
\setlength\baselineskip{2cm}
\renewcommand{\baselinestretch}{1.2}

\makeindex

\title{
{\Large\sf Relat�rio Descritivo} \vspace{0.5cm} \\
\textbf{\LARGE Sistema de localiza��o por fus�o GPS/INS/Alt�mero}\\
}
\author{Geovany Ara�jo Borges \vspace{2mm}\\
\it Laborat�rio de Rob�tica e Automa��o (LARA) \\
\it Grupo de Automa��o, Rob�tica e Vis�o Computacional (GRAV)\\
\it Departamento de Engenharia El\'{e}trica - FT - UnB\\
\it Caixa Postal 4.386 - Bras\'{i}lia - DF - Brasil\\
}

\begin{document}
\maketitle

\begin{abstract}
O resumo deve ser iniciado por uma breve frase expondo o principal objeto do relat�rio. Em seguida suas caracter�sticas s�o mostradas. Simula��es e experimentos, se houver, devem ser brevemente comentados.
\end{abstract}

\section{Introdu��o} 

\indent Um primeiro bl� introduzindo ao leitor o contexto do problema, uma breve descri��o, o que foi realizado e o teor do resto do documento \cite{book:johnson:1980}.

Esse modelo est� pronto para ser aberto no Texniccenter, bastando clicar em \texttt{report.tcp}. Enquano estiver editando, visualize rapidamente o resultado compilando para gerar apenas \texttt{.DVI} (op��o \texttt{Latex => DVI}). Para a vers�o final, gere \texttt{.PDF} com a op��o \texttt{Latex => PS => PDF}. As figuras s�o em formato \texttt{.EPS}. Gere a figura em qualquer editor gr�fico e converta para um formato conhecido (\texttt{.PNG}, \texttt{.JPG}, etc.). Abra a figura no GIMP e converta para \texttt{.EPS}.

O restante deste documento est� organizado nas seguintes se��es. A Se��o \ref{sec:secprimeira} descreve uma coisa. Uma outra coisa � descrita na Se��o \ref{sec:secsegunda}. Considerando-se o que foi apresentado, a Se��o \ref{sec:conclusoes} apresenta as conclus�es deste relat�rio. 

\section{Se��o primeira} \label{sec:secprimeira}

\section{Se��o segunda} \label{sec:secsegunda}

\begin{figure}
\centering
\includegraphics[width=10cm]{figs/unb.eps}
\caption{Figura exemplo.}
\label{fig:devhelp}
\end{figure} 

\section{Conclus�es}\label{sec:conclusoes}

\bibliographystyle{unsrt} 
\bibliography{report}

\appendix

\end{document}          

